% Source: Wald GR, physical PDF pp. 447--449
%         (printed pp. 434--436).
% Coverage: Appendix B.3, Frobenius's Theorem; equations (B.3.1)--(B.3.6);
%           Theorems B.3.1--B.3.2.
% Figures/tables/footnotes: none.
% Status: source-reviewed and compile-clean.
% Uncertainties: none.

\subsection{Frobenius's Theorem}\label{sec:B.3}

Let \(M\) be an \(n\)-dimensional manifold. An issue which arises frequently is the
following: At each point \(x\in M\) we are given a subspace
\(W_x\subsetneq T_xM\) of the tangent space \(T_xM\) with
\(\dim W_x=m<n\). The subspace \(W_x\) is required to vary smoothly with \(x\) in
the sense that for each \(x\in M\) we can find an open neighborhood \(O\) of \(x\)
such that in \(O\), \(W\) is spanned by \(C^\infty\) vector fields. We denote the
collection of subspaces \(W_x\) by \(W\). We wish to know whether we can find
integral submanifolds of \(W\), i.e., whether through each point \(x\) we can find an
embedded submanifold \(S\) such that the tangent space \(T_yS\) to this submanifold
at each \(y\in S\) coincides with \(W_y\). An important special case of this general
problem arises when one has a metric on \(M\) and wishes to know if a vector field
\(\xi^a\) is orthogonal to a family of hypersurfaces (see, e.g., Section~6.1), i.e.,
whether the \((n-1)\)-dimensional subspaces, \(W\), orthogonal to \(\xi^a\) are
integrable.

If the subspaces \(W\) are one-dimensional, the above problem reduces to that of
finding integral curves of a smooth vector field \(v^a\). As discussed in
Section~2.2, such integral curves always can be found. However, if
\(\dim W>1\), it is possible for the \(W\)-planes to \lq\lq twist around\rq\rq\ so
that integral submanifolds cannot be found. To see that this is the case, we note that
if we could find integral submanifolds, we could span \(W\) in a neighborhood of any
point by coordinate vector fields \(X_1^a,\ldots,X_m^a\) in \(M\) such that
\([X_\mu,X_\nu]=0\). Since any two vector fields \(Y^a,Z^a\) which lie in \(W\) can
be expressed as linear combinations of these coordinate vector fields, this implies
that for all \(Y^a,Z^a\in W\) we have
\begin{equation}
  [Y,Z]^a
  =\sum_{\mu,\nu}[f_\mu X_\mu,g_\nu X_\nu]^a
  =\sum_{\mu,\nu}
  \bigl(f_\mu X_\mu(g_\nu)-g_\mu X_\mu(f_\nu)\bigr)X_\nu^a\in W.
  \tag{B.3.1}\label{eq:B.3.1}
\end{equation}
If \(W\) satisfies the property that \([Y,Z]^a\in W\) for all
\(Y^a,Z^a\in W\), then \(W\) is said to be \emph{involute}. We have just shown that
a necessary condition for \(W\) to possess integral submanifolds is that it be
involute. Conversely, it can be shown (see, e.g., Bishop and Crittenden 1964) that
this condition is also sufficient. This result is known as Frobenius's theorem.

\renewcommand{\thetheorem}{B.3.\arabic{theorem}}
\setcounter{theorem}{0}
\begin{theorem}[Frobenius's theorem; vector form]\label{thm:B.3.1}
A necessary and sufficient condition for a smooth specification, \(W\), of
\(m\)-dimensional subspaces of the tangent space at each \(x\in M\) to possess
integral submanifolds is that \(W\) be involute, i.e., for all
\(Y^a,Z^a\in W\) we have \([Y,Z]^a\in W\).
\end{theorem}

Frobenius's theorem also has a dual formulation in terms of differential forms.
Given \(W_x\subsetneq T_xM\) as above, we can consider the one-forms
\(\omega\in T_x^\vee M\) which satisfy
\begin{equation}
  \omega_aX^a=0
  \tag{B.3.2}\label{eq:B.3.2}
\end{equation}
for all \(X^a\in W_x\). It is not difficult to see that such \(\omega\)'s span an
\((n-m)\)-dimensional subspace,
\(\mathcal T_x^\vee\subseteq T_x^\vee M\), of the cotangent space at \(x\). Conversely, an
\((n-m)\)-dimensional subspace \(\mathcal T_x^\vee\) of \(T_x^\vee M\) defines an
\(m\)-dimensional subspace \(W_x\) of \(T_xM\) via Eq.~\eqref{eq:B.3.2}. Thus, we may
reformulate our above question in terms of \(\mathcal T^\vee\): Under what conditions
does a smooth specification, \(\mathcal T^\vee\), of \((n-m)\)-dimensional subspaces
of one-forms at each point have the property that the associated tangent subspaces
\(W\) (consisting at each \(x\) of all vectors \(X^a\) satisfying
\(\omega_aX^a=0\) for all \(\omega_a\in\mathcal T_x^\vee\)) admit integral
submanifolds?

According to Frobenius's theorem, integral submanifolds will exist if and only if for
all \(\omega_a\in\mathcal T^\vee\) and all \(Y^a,Z^a\in W\) (so that
\(\omega_aY^a=\omega_aZ^a=0\)), we have
\begin{equation}
  \omega_a[Y,Z]^a=0.
  \tag{B.3.3}\label{eq:B.3.3}
\end{equation}
To see what this implies for \(\omega_a\), we substitute our expression
\eqref{eq:3.1.2} for the commutator in terms of an arbitrary derivative operator
\(\nabla_a\) to obtain
\begin{equation}
  \begin{aligned}
  0
  & =\omega_a(Y^b\nabla_bZ^a-Z^b\nabla_bY^a)\\
  & =-Z^aY^b\nabla_b\omega_a+Y^aZ^b\nabla_b\omega_a\\
  & =2Y^aZ^b\nabla_{[b}\omega_{a]}.
  \end{aligned}
  \tag{B.3.4}\label{eq:B.3.4}
\end{equation}
However, Eq.~\eqref{eq:B.3.4} can hold for \(Y^a\) and \(Z^a\) in the subspace
annihilated by \(\mathcal T^\vee\) if and only if
\(\nabla_{[a}\omega_{b]}\) can be expressed as
\begin{equation}
  \nabla_{[a}\omega_{b]}
  =\sum_{\alpha=1}^{n-m}\mu^\alpha_{[a}\nu^\alpha_{b]},
  \tag{B.3.5}\label{eq:B.3.5}
\end{equation}
where each \(\nu^\alpha\) is an arbitrary one-form and each
\(\mu^\alpha\in\mathcal T^\vee\). Thus, we can reformulate Frobenius's theorem in terms
of differential forms as follows:

\begin{theorem}[Frobenius's theorem; dual formulation]\label{thm:B.3.2}
Let \(\mathcal T^\vee\) be a smooth specification of an \((n-m)\)-dimensional subspace
of one-forms. Then the associated \(m\)-dimensional subspace \(W\) of the tangent
space admits integral submanifolds if and only if for all
\(\boldsymbol\omega\in\mathcal T^\vee\) we have
\[
  \dd\boldsymbol\omega
  =\sum_\alpha\boldsymbol\mu^\alpha\wedge\boldsymbol\nu^\alpha,
\]
where each \(\boldsymbol\mu^\alpha\in\mathcal T^\vee\).
\end{theorem}

The dual formulation of Frobenius's theorem gives a useful criterion for when a
vector field \(\xi^a\) is hypersurface orthogonal. Letting \(\mathcal T^\vee\) be the
one-dimensional subspace spanned by \(\xi_a=g_{ab}\xi^b\), we see that \(\xi^a\)
will be hypersurface orthogonal if and only if
\(\nabla_{[a}\xi_{b]}=\xi_{[a}\nu_{b]}\) (where we have set
\(\mu_a=\xi_a\) since \(\mathcal T^\vee\) is one-dimensional). This latter condition is
equivalent to \(\xi_{[a}\nabla_b\xi_{c]}=0\), and thus we see that the necessary and
sufficient condition that \(\xi^a\) be hypersurface orthogonal is
\begin{equation}
  \xi_{[a}\nabla_b\xi_{c]}=0.
  \tag{B.3.6}\label{eq:B.3.6}
\end{equation}
